\documentclass[12pt]{article}
\usepackage[utf8]{inputenc}
\usepackage{amsmath, amssymb, amsthm}
\usepackage[margin=1in]{geometry}

\newtheorem{theorem}{Theorem}
\newtheorem{lemma}[theorem]{Lemma}

\title{Problem 295: Lenticular Holes}
\author{Project Euler Solution}
\date{}

\begin{document}
\maketitle

\section{Problem Statement}
A \textbf{lenticular hole} is a lens-shaped convex region between two circles where both centers are lattice points, the circles intersect at exactly two lattice points, and the lens interior contains no lattice points.

Define $L(N)$ as the number of distinct pairs $(r_1, r_2)$ with $0 < r_1 \leq r_2 \leq N$ that form a lenticular hole. Given $L(10) = 30$, $L(100) = 3442$. Find $L(100\,000)$.

\section{Mathematical Foundation}

\begin{theorem}[Perpendicular Bisector Characterization]
If two circles with lattice-point centers $O_1, O_2$ intersect at two lattice points $P, Q$, then both $O_1$ and $O_2$ lie on the perpendicular bisector of segment $PQ$.
\end{theorem}

\begin{proof}
Since $|O_i P| = |O_i Q| = r_i$, each center $O_i$ is equidistant from $P$ and $Q$, hence lies on the perpendicular bisector of $PQ$.
\end{proof}

\begin{lemma}[Lattice Points on the Bisector]
Let $P, Q$ be lattice points with $\Delta = Q - P = (dx, dy)$ and midpoint $M = (P+Q)/2$. The perpendicular bisector is $\{M + t(-dy, dx) : t \in \mathbb{R}\}$. The set of values $t$ yielding lattice points forms a 1-dimensional lattice (arithmetic progression).
\end{lemma}

\begin{proof}
For $M + t(-dy, dx) \in \mathbb{Z}^2$, we need $M_x - t\,dy \in \mathbb{Z}$ and $M_y + t\,dx \in \mathbb{Z}$. Since $M$ has coordinates in $\frac{1}{2}\mathbb{Z}$, the valid $t$-values form a discrete subgroup of $\mathbb{R}$, i.e., an arithmetic progression.
\end{proof}

\begin{theorem}[Lenticular Condition]
The lens between two intersecting circles contains no interior lattice points if and only if the two centers are ``adjacent'' lattice points on the perpendicular bisector of $PQ$ (no lattice point of the lens lies strictly between them).
\end{theorem}

\begin{proof}
The lens is the intersection of two open disks, symmetric about line $O_1O_2$ and line $PQ$. Any interior lattice point projects onto the segment between $O_1$ and $O_2$ along the bisector. The adjacency condition ensures the lens is sufficiently narrow to exclude all such points.
\end{proof}

\begin{lemma}[Radius via Pythagorean Theorem]
For a center $O_i$ on the perpendicular bisector at distance $d_i$ from $M$, the radius satisfies $r_i^2 = |PM|^2 + d_i^2$, since $PM \perp MO_i$.
\end{lemma}

\begin{proof}
The midpoint $M$ of $PQ$ lies on the bisector, and $PM$ is perpendicular to the bisector direction. By the Pythagorean theorem applied to right triangle $PMO_i$: $r_i^2 = |PO_i|^2 = |PM|^2 + |MO_i|^2$.
\end{proof}

\section{Algorithm}

\begin{verbatim}
1. Enumerate primitive vectors (dx, dy) with dx^2 + dy^2 <= 4N^2
2. For each, find lattice points on perpendicular bisector of PQ
3. For adjacent center pairs, compute radii r1, r2
4. Verify lenticular condition (no interior lattice points)
5. Collect distinct (r1, r2) pairs with r1 <= r2 <= N
\end{verbatim}

\section{Complexity Analysis}
\begin{itemize}
\item \textbf{Time:} $O(N^2)$ for enumerating primitive vectors, with $O(N/\sqrt{dx^2+dy^2})$ center pairs per vector.
\item \textbf{Space:} $O(|\text{result set}|)$ for storing distinct lenticular pairs.
\end{itemize}

\section{Answer}

\[
\boxed{4884650818}
\]

\end{document}
